% Generated draft: requires mathematical and source review.
\documentclass[11pt]{article}
\usepackage[margin=25mm]{geometry}
\usepackage{fontspec}
\setmainfont{Latin Modern Roman}
\usepackage{amsmath,amssymb,mathtools,unicode-math}
\setmathfont{Latin Modern Math}
\usepackage{enumitem,needspace,xurl,hyperref}
\setlength{\parindent}{0pt}
\setlength{\parskip}{5pt}
\setlength{\emergencystretch}{3em}
\newcommand{\N}{\mathbb{N}}
\newcommand{\Z}{\mathbb{Z}}
\newcommand{\Q}{\mathbb{Q}}
\newcommand{\R}{\mathbb{R}}
\newcommand{\C}{\mathbb{C}}
\newcommand{\F}{\mathbb{F}}
\newcommand{\A}{\mathbb{A}}
\newcommand{\PP}{\mathbb{P}}
\newcommand{\E}{\mathbb{E}}
\newcommand{\eps}{\varepsilon}
\newcommand{\dd}{\,\mathrm{d}}
\begin{document}
\section*{UnsolvedMath ID: 800233}
\textbf{Source problem code:} catalogue entry
\subsection{Definitions and mathematical statement}
Let \(n\geq 1\) be an integer and let \(\Omega\subset\C^n\) be a bounded domain. Interpret the source's phrase ``real analytic domain'' as requiring real analytic boundary: near each boundary point there is a real analytic defining function \(\rho\) such that \(\Omega=\{\rho<0\}\) locally and \(\dd\rho\ne0\) on the boundary. Assume that \(\Omega\) is pseudoconvex, meaning that it admits a continuous plurisubharmonic exhaustion function. The boundary interpretation is conventional but is not explicitly specified in the supplied statement.

For \(0<\beta<1\), a function \(v\) belongs to \(C^{0,\beta}(\overline\Omega)\) if
\[
\|v\|_{C^{0,\beta}(\overline\Omega)}:=\sup_{z\in\overline\Omega}|v(z)|+\sup_{\substack{z,w\in\overline\Omega\\z\ne w}}\frac{|v(z)-v(w)|}{|z-w|^\beta}<\infty.
\]
For differential forms, use the sum of these norms over the coefficients in standard coordinates. Write \(\bar\partial\) for the Dolbeault operator.

\textbf{Source question (Problem 30).} Can one solve \(\bar\partial\) with H\"older estimates on every bounded pseudoconvex domain in \(\C^n\) with real analytic boundary?

To describe the unresolved meaning of ``H\"older estimates,'' one possible estimate template, after specifying a form degree \(1\leq q\leq n\), a normed data space \(X^{0,q}(\Omega)\), and an exponent \(0<\beta<1\), is
\[
\forall f\in X^{0,q}(\Omega),\qquad
\bar\partial f=0\ \Longrightarrow\ \exists u\in C^{0,\beta}_{(0,q-1)}(\overline\Omega):\quad
\bar\partial u=f,\qquad
\|u\|_{C^{0,\beta}(\overline\Omega)}\leq C\|f\|_{X^{0,q}(\Omega)}.
\]
Here the equations may be understood distributionally. This display explains a possible type of estimate; it is not an assertion that these particular spaces or quantifiers are intended by the source. The supplied question does not specify the form degree, data regularity, H\"older exponent or gain, endpoint conventions, or permitted dependence of the estimate constant. It therefore does not determine a unique quantified estimate. In particular, the source does not explicitly say whether the estimates are global up to the boundary, as in this template, or local.
\subsection{Short English statement}
The question asks whether every bounded pseudoconvex domain with real analytic boundary admits solutions of \(\bar\partial u=f\) with H\"older control. The supplied wording leaves the data spaces, form degrees, exponents, and precise estimates unspecified; these require mathematical clarification.
\subsection{Sources}
\begin{itemize}
\item Supplied local UnsolvedMath registry record, permanent ID \(800233\), source code \texttt{catalogue entry}, titled ``Several Complex Variables Problem List --- Problem 30.''
\item Buzzard, Fornaess, Krantz, \emph{The Several Complex Variables Problem List} (1995), Problem 30. Supplied source URL: \url{https://arxiv.org/abs/math/9509202}. The registry attributes the problem to Buzzard, Fornaess, Krantz, with proposed year 1995.
\item Supplied provenance records access on 2026-10-16 and extraction method \texttt{arxiv\_latex}. These are supplied provenance claims, not an independent source visit or verification. The registry labels the problem ``open'' and records an update on 2026-10-16; this status is supplied curation, not an independently verified mathematical or literature assessment.
\end{itemize}
\end{document}
